\chapter{Механические свойства и разрушение}\label{ch:mech}

Ради чего вообще считают гидриды? Ради того, чтобы знать, выдержит ли оболочка. Эта глава --- о том, как
механические свойства измеряют, что означают числа в паспорте сплава и как гидриды превращают пластичный металл в
хрупкий.

\section{Испытание на растяжение}

Образец растягивают с постоянной скоростью и записывают силу $F$ и удлинение $\Delta L$. Из них получают
\textbf{условные} (инженерные) напряжение и деформацию, отнесённые к исходным размерам:
\begin{equation}
s = \frac{F}{A_0},\qquad e = \frac{\Delta L}{L_0},
\end{equation}
и \textbf{истинные}, отнесённые к текущим:
\begin{equation}
\sigma = \frac{F}{A} = s(1+e),\qquad \eps = \int\frac{\dd L}{L} = \ln(1+e).
\end{equation}
Последние формулы используют сохранение объёма при пластической деформации, $AL = A_0L_0$, и верны до образования
шейки. Истинная деформация аддитивна: два последовательных растяжения на $\eps_1$ и $\eps_2$ дают $\eps_1 + \eps_2$.

\fig{k1_tensile}{Модельный сплав циркония при комнатной температуре: $E = 95$~ГПа, упрочнение по Холломону
$\sigma = K\eps_p^n$ ($K = 750$~МПа, $n = 0.10$). а) Условная (сплошная) и истинная (пунктир) диаграммы; предел
прочности $\sigma_{\text{в}}$ --- начало шейки. б) Условный предел текучести $\sigma_{0.2}$: пересечение кривой с прямой,
параллельной упругому участку и сдвинутой на 0.2\,\%. в) Критерий Консидера: шейка начинается, когда скорость
упрочнения $\dd\sigma/\dd\eps$ падает до самого напряжения; для закона Холломона --- при $\eps_u = n$.}

Характеристики, которые выдаёт испытание (рис.~\ref{fig:k1_tensile}):
\begin{itemize}
\item \textbf{модуль упругости} $E$ --- наклон начального участка (его точнее мерят динамически, по скорости звука или
  резонансу);
\item \textbf{условный предел текучести} $\sigma_{0.2}$ --- напряжение, при котором остаточная деформация 0.2\,\%. У
  большинства металлов нет резкого перехода от упругости к пластичности, и порог приходится задавать условно;
\item \textbf{предел прочности} (временное сопротивление) $\sigma_{\text{в}}$ --- максимум условного напряжения;
\item \textbf{равномерное удлинение} --- деформация до шейки, и \textbf{полное удлинение} --- до разрыва; \textbf{сужение}
  площади в шейке.
\end{itemize}

\textbf{Почему появляется шейка.} Сила $F = \sigma A$. При растяжении $\sigma$ растёт за счёт упрочнения, а $A$ падает.
Пока упрочнение побеждает, любое случайное утонение упрочняется сильнее соседей и выравнивается. Когда $\dd F = 0$,
то есть $\dd\sigma/\sigma = -\dd A/A = \dd\eps$, равновесие становится неустойчивым, и деформация локализуется:
\begin{keyf}[Критерий Консидера]
\[
\frac{\dd\sigma}{\dd\eps} = \sigma\quad\Longrightarrow\quad \eps_u = n \text{ для }\sigma = K\eps^n .
\]
\end{keyf}
Показатель упрочнения $n$ прямо равен равномерной деформации. Холоднодеформированный металл упрочняется слабо ($n$
мал), поэтому прочнее, но шейка в нём наступает раньше.

\begin{zrbox}
Оболочку нагружают давлением, и важна \emph{окружная} прочность и пластичность. Её измеряют на кольцах:
\begin{itemize}
\item \textbf{растяжение кольца} на двух полудисках (D-образных оправках) --- окружное направление растягивается, но у
  оправок есть трение и изгиб, поэтому кривая не чисто одноосная, и деформацию обычно относят к рабочим участкам;
\item \textbf{сжатие кольца} между плитами --- простой способ оценить охрупчивание гидридами: кольцо сплющивают и
  смотрят, при каком смещении появляется трещина; при радиальных гидридах она возникает гораздо раньше;
\item \textbf{испытание трубы внутренним давлением} --- ближе всего к реальности (двухосное нагружение
  $\sigma_z = \sigma_\theta/2$ для трубы с заглушками).
\end{itemize}
Из-за текстуры свойства в окружном, осевом и радиальном направлениях различаются. Это описывают коэффициентами
анизотропии (Ланкфорда, Хилла; глава~\ref{ch:plasticity}).
\end{zrbox}

\section{Твёрдость}

Твёрдость --- сопротивление вдавливанию. По Виккерсу в металл вдавливают алмазную пирамидку силой $F$ и меряют
диагональ отпечатка $d$:
\begin{equation}
HV = 1.854\,\frac{F}{d^2}
\end{equation}
($F$ в кгс, $d$ в мм --- исторические единицы). Под индентором металл течёт, и твёрдость связана с пределом
текучести: $HV\approx3\sigma_y$ в согласованных единицах (для металлов с малым упрочнением). Микротвёрдость меряют
с нагрузками в граммы --- на отдельных фазах; гидрид заметно твёрже матрицы. Наноиндентирование даёт и модуль, и
твёрдость в объёмах меньше микрометра.

\section{Ударная вязкость и переход хрупкость--пластичность}

Испытание по Шарпи: образец с надрезом ломают ударом маятника и меряют поглощённую энергию. ГЦК металлы пластичны при
любой температуре. ОЦК металлы (сталь) и многие ГПУ с понижением температуры переходят от вязкого излома к
хрупкому в узком интервале --- \textbf{температура хрупко-вязкого перехода}. Причина: предел текучести растёт при
охлаждении быстрее, чем сопротивление скола, и при некоторой температуре скол наступает раньше течения.

\begin{zrbox}
Циркониевые сплавы без гидридов пластичны и при комнатной температуре. Гидриды создают в них переход
хрупкость--пластичность: при комнатной температуре гидрид ломается почти без пластичности, а при 250--350\degC{}
он заметно пластичнее и матрица вокруг него легче течёт. Радиальные гидриды сдвигают переход вверх по температуре
на десятки--сотни градусов: трещина идёт по ним сквозь стенку. Именно поэтому критерием безопасности при
хранении отработавшего топлива служит не количество водорода, а доля и связность радиальных гидридов.
\end{zrbox}

\section{Механика разрушения}

\textbf{Почему реальная прочность мала.} Теоретическая прочность связи --- порядка $E/10$ (упражнение~\ref{ch:atoms}.1),
в десятки--сотни раз больше реальной. Гриффитс понял, что всё дело в трещинах: у вершины трещины напряжение концентрируется.
Трещина длины $2a$ в пластине под напряжением $\sigma$ растёт, если высвобождаемая упругая энергия покрывает работу
создания новых поверхностей. Для хрупкого материала
\begin{equation}
\sigma_c = \sqrt{\frac{2E\gamma_s}{\pi a}},
\end{equation}
$\gamma_s$ --- поверхностная энергия. Критическое напряжение падает как $1/\sqrt a$.

\textbf{Коэффициент интенсивности напряжений.} Около вершины любой трещины в упругом теле напряжения одинаково
устроены:
\begin{equation}
\sigma_{yy}(r) = \frac{K}{\sqrt{2\pi r}},\qquad K = Y\sigma\sqrt{\pi a},
\end{equation}
где $r$ --- расстояние от вершины, а $Y$ --- безразмерный множитель формы (1 для центральной трещины в бесконечной
пластине, 1.12 для краевой). Вся информация о нагрузке и размере трещины свёрнута в одно число $K$. Трещина растёт,
когда $K$ достигает критического значения --- \textbf{вязкости разрушения} $K_c$ (при плоской деформации --- $K_{Ic}$),
свойства материала.
\begin{keyf}[Условие разрушения]
\[
K = Y\sigma\sqrt{\pi a} = K_{c}.
\]
\end{keyf}

\fig{k2_fracture}{а) Напряжение перед вершиной трещины при $K = 20$~МПа$\sqrt{\text{м}}$ и $\sigma_y = 500$~МПа: упругое
решение $K/\sqrt{2\pi r}$ и его срез на пределе текучести; размер пластической зоны по Ирвину. б) Критическое
напряжение для краевой трещины глубины $a$ при разной вязкости разрушения (оценка, $Y = 1.12$); пунктир --- толщина
стенки оболочки 0.6~мм.}

В металле напряжение у вершины не может превысить предел текучести: перед трещиной образуется \textbf{пластическая зона}
размером
\begin{equation}
r_p\approx\frac{1}{2\pi}\Bigl(\frac{K}{\sigma_y}\Bigr)^2,
\end{equation}
и на её создание уходит гораздо больше энергии, чем на поверхность. Поэтому $K_c$ металлов --- десятки
МПа$\sqrt{\text{м}}$, а керамик и гидридов --- единицы. Линейная механика разрушения применима, пока $r_p$ мала по
сравнению с трещиной и сечением; иначе используют $J$-интеграл и критерии раскрытия трещины.

\begin{zrbox}
Радиальный гидрид в стенке --- готовая трещина: он почти не сопротивляется разрушению ($K_c$ гидрида --- единицы
МПа$\sqrt{\text{м}}$), и под окружным напряжением его вершина работает как вершина трещины в матрице. Цепочка
радиальных гидридов длиной в сотни микрометров при вязкости матрицы с гидридами в десяток МПа$\sqrt{\text{м}}$
даёт критическое напряжение в сотни МПа (рис.~\ref{fig:k2_fracture}б) --- того же порядка, что напряжения в оболочке.
Окружный гидрид той же длины лежит вдоль напряжения и такой опасности не создаёт.
\end{zrbox}

\section{Замедленное гидридное растрескивание}

У вершины трещины гидростатическое растяжение велико, и водород стекается туда (глава~\ref{ch:diffusion}):
химический потенциал водорода ниже там, где решётка растянута. У вершины концентрация превышает предел выпадения, и
там вырастает гидрид, ориентированный поперёк нагрузки. Когда он достигает критического размера, он ломается,
трещина продвигается на его длину --- и цикл повторяется. Это \textbf{замедленное гидридное растрескивание} (ЗГР, DHC).
У него есть порог $K_{IH}$ --- ниже него трещина не растёт --- и скорость, зависящая от температуры через диффузию и
растворимость. Существенно, что ЗГР может идти при водороде ниже предела растворимости при данной температуре:
водород приходит к вершине из окружающего раствора.

\section{Ползучесть и релаксация}

При температурах выше примерно трети температуры плавления (для циркония --- выше 300--400\degC) металл медленно
деформируется при постоянном напряжении ниже предела текучести --- \textbf{ползёт}. Кривая ползучести
(рис.~\ref{fig:k3_creep}а) имеет три стадии: неустановившуюся (скорость падает: упрочнение побеждает возврат),
установившуюся (упрочнение и возврат уравновешены) и ускоренную (повреждения, утонение, разрушение). Скорость
установившейся ползучести описывают степенным законом
\begin{equation}
\dot\eps = A\,\sigma^n\exp\Bigl(-\frac{Q}{RT}\Bigr).
\end{equation}
Показатель $n\approx3$--7 и энергия активации, близкая к энергии самодиффузии, указывают на дислокационную ползучесть
(переползание дислокаций). При малых напряжениях $n\to1$ --- диффузионная ползучесть: атомы перетекают от сжатых
границ зёрен к растянутым (механизмы Набарро--Херринга и Кобла).

\fig{k3_creep}{а) Кривая ползучести при постоянной нагрузке: неустановившаяся (I), установившаяся (II) и
ускоренная (III) стадии, крест --- разрушение. б) Релаксация напряжения при постоянной полной деформации для
степенного закона $\dot\eps = A\sigma^5$ при трёх значениях $A$ (условно --- трёх температурах): чем горячее, тем
быстрее спадает напряжение.}

\textbf{Релаксация} --- обратная ситуация: полная деформация зафиксирована, а напряжение спадает, потому что упругая
деформация постепенно превращается в ползучесть. Из $\eps = \sigma/E + \eps_{\text{полз}} = \text{const}$ следует
$\dot\sigma = -E\dot\eps_{\text{полз}} = -EA\sigma^n$, откуда
\begin{equation}
\sigma(t) = \sigma_0\Bigl[1 + (n-1)EA\sigma_0^{\,n-1}t\Bigr]^{-1/(n-1)}
\end{equation}
(рис.~\ref{fig:k3_creep}б). Напряжение падает сначала быстро, затем всё медленнее.

\begin{zrbox}
Ползучесть и релаксация важны для переориентации по двум причинам. Во-первых, при сушке и хранении отработавшего
топлива окружное напряжение от газа под оболочкой медленно уменьшается --- за счёт охлаждения и ползучести, --- и
выпадение гидридов идёт при меняющемся напряжении. Во-вторых, в опытах на кольцах важно, как задана нагрузка: при
постоянной силе напряжение держится, при постоянном смещении --- релаксирует во время выдержки и охлаждения, и
гидриды выпадают уже при меньшем напряжении, чем было приложено вначале.
\end{zrbox}

\begin{exercise}
Докажите, что для закона Холломона $\sigma = K\eps^n$ шейка начинается при $\eps = n$. Чему тогда равен предел прочности
$\sigma_{\text{в}}$ (условное напряжение)?
\end{exercise}

\begin{exercise}
Оцените глубину краевой трещины, при которой оболочка с $K_c = 15$~МПа$\sqrt{\text{м}}$ разрушится при окружном
напряжении 300~МПа. Сравните с толщиной стенки.
\end{exercise}

\begin{exercise}
Выведите формулу релаксации напряжения для степенного закона ползучести. За какое время напряжение упадёт вдвое?
\end{exercise}
